% Section 1.3: 研究问题与挑战

尽管大语言模型在代码生成任务上已取得长足进展，但在面向算法题的场景下，从“可用”走向“稳定可用”仍面临若干关键问题与挑战。结合现有研究现状与本文的实验观察，可以将这些问题归纳为以下四个方面。

\paragraph{（1）推理过程不稳定，可执行正确性波动较大。}
算法代码生成本质上依赖严谨的逻辑推导与边界条件处理，而通用大模型在未经针对性适配时，常表现为“同题不同解”的高方差现象：同一问题在不同采样下可能产生逻辑错误、循环越界或遗漏特殊用例的实现。这种不稳定性使得模型即使在统计平均上表现尚可，也难以满足工程场景对确定性的要求。

\paragraph{（2）复杂问题上的规划能力不足。}
对于涉及动态规划、图搜索或数学推导的问题，模型需要在生成代码之前完成隐式的算法规划。然而，直接以自回归方式生成代码往往使模型“跳过”规划环节，导致解法在整体结构上偏离最优思路，或将本应在外层处理的状态维护内嵌到循环内部，造成逻辑耦合与错误扩散。如何显式引导模型先行完成问题抽象与步骤规划，是提升复杂算法题表现的关键。

\paragraph{（3）训练资源受限下的能力迁移问题。}
对数十亿参数规模的基座模型进行全量微调，对单机训练环境而言代价高昂；而直接使用通用指令微调模型在算法题场景下又难以充分发挥数据价值。因此，如何在有限算力下，借助参数高效微调方法在目标任务上稳定地获取增益，并保证训练过程的可复现性，是本文必须解决的工程性问题。

\paragraph{（4）推理引导与训练方法的协同设计问题。}
LoRA 在训练侧提供的是参数层面的任务适配，CoT 在推理侧提供的是结构层面的推理引导，二者属于不同技术栈。然而，它们对模型在算法题上的最终表现是否具有正交收益、Few-shot 示例在 LoRA 模型与基座模型上是否展现出一致的边际效应，现有工作缺乏统一的对照实验。如何在同一实验协议下清晰地刻画“Base / LoRA / Zero-shot CoT / Few-shot CoT”的相对优劣，是本文在实验设计层面需要重点回答的问题。

围绕上述四类问题，本文将在后续章节通过数据构建、LoRA 微调、双 CoT 对照以及统一的可执行评测流程，系统地给出针对性的解决方案与实证分析。